Put \(E=e^\eta\), \(k=\lfloor(d-1)/2\rfloor\), and choose
\[
\alpha=\frac{E-D}{8(E-1)}\in(0,1/8),
\qquad p=\frac{\alpha}{D}.
\]
Take an orthonormal system \(e_0,(u_j,w_j)_{j=1}^k\), completed by a vector \(z\) when \(d\) is even, and use actions \(\{+,-,0\}\). At each hard context \(x_j\), set
\[
\pi_{\mathrm{ref}}(+\mid x_j)=\pi_{\mathrm{ref}}(-\mid x_j)=p,\qquad
\pi_{\mathrm{ref}}(0\mid x_j)=1-2p,
\]
\[
\phi(x_j,+)=\frac{u_j+w_j}{\sqrt2},\qquad
\phi(x_j,-)=\frac{u_j-w_j}{\sqrt2},\qquad
\phi(x_j,0)=0.
\]
Give each hard context raw mass \(1/k\). To calibrate the \(e_0\) block, add a context \(x_C\) of raw mass one, put
\[
q=\frac{E/C-1}{E-1},\qquad
\pi_{\mathrm{ref}}(+\mid x_C)=q,\qquad
\phi(x_C,+)=e_0,\qquad
\phi(x_C,-)=\phi(x_C,0)=0,
\]
and distribute the remaining reference mass positively over \(-,0\). If \(C>D\), add a context \(x_A\) of raw mass
\[
a_{\mathrm{anc}}=q\frac{C-D}{D-1}
\]
at which every action feature equals \(e_0\) and give all actions positive reference probability; if \(C=D\), omit it. If \(z\) is present, add a context of raw mass one at which every action feature equals \(z\), and give all actions positive reference probability. Normalize all raw masses, and denote their sum by \(H\). This sum is bounded above and below by positive constants depending only on the coverage and regularization indices, uniformly in the dimension. This public experiment is fixed once \((d,C,D,\eta)\) is fixed.

For \(t=\eta\gamma\), define
\[
T_\gamma=\frac{D(1-2p)}{e^t-2\alpha\cosh t},
\qquad \beta_\gamma=\eta^{-1}\log T_\gamma.
\]
At \(\gamma=0\),
\[
T_0=\frac{D-2\alpha}{1-2\alpha}.
\]
This number is strictly larger than one. It is also strictly smaller than \(E\), because \(2\alpha(E-1)=(E-D)/4<E-D\). The denominator, \(T_\gamma\), and \(\beta_\gamma\) are continuous near zero. Therefore there are \(\gamma_0>0\) and \(m\in(0,1/2)\), depending only on \((D,\eta)\), such that for \(0\le\gamma\le\gamma_0\),
\[
t\le1,\qquad m\le\beta_\gamma-\gamma\le\beta_\gamma+\gamma\le1-m,
\]
and the denominator defining \(T_\gamma\) is positive.

For each \(v\in\{-1,1\}^k\), set the coefficient in direction \(e_0\) equal to one, set
\[
\theta_{u_j}=\sqrt2\beta_\gamma,\qquad
\theta_{w_j}=\sqrt2\gamma v_j,
\]
and set the optional \(z\)-coefficient to zero. Thus the hard-context rewards are \(\beta_\gamma\pm\gamma v_j\), while the rewards at \(x_C\) are one for action \(+\) and zero for the other actions. Use Bernoulli rewards with these means.

We verify exact-shell membership. At \(x_C\), the action \(+\) target-to-reference ratio is exactly \(C\), and every other ratio is \(C/E<1\). At a hard context the Gibbs normalizer is
\[
Z=1-2p+2pT_\gamma\cosh t.
\]
The definition of \(T_\gamma\) is equivalent to \(T_\gamma e^t=DZ\). Hence the favored and unfavored informative actions have target-to-reference ratios \(D\) and \(De^{-2t}\). Since \(\beta_\gamma\ge\gamma\), one has \(T_\gamma e^t\ge1\), so the zero-action ratio \(1/Z\) is at most \(D\). All ratios at action-invariant contexts equal one. Because \(D\le C\), this proves \(C_P=C\).

In the \((u_j,w_j)\) basis, the logging covariance block is a positive context-mass multiple of \(pI_2\). The target probabilities of the favored and unfavored actions are respectively
\[
pD=\alpha,\qquad pDe^{-2t}=\alpha e^{-2t},
\]
so the two generalized eigenvalues of the target block are \(D\) and \(De^{-2t}\). In the \(e_0\) direction, the generalized ratio is \(D\), because it is \(C=D\) if the anchor is absent and otherwise
\[
\frac{qC+a_{\mathrm{anc}}}{q+a_{\mathrm{anc}}}=D.
\]
The optional \(z\) ratio is one. Every basis direction has positive logging mass, so \(\Lambda_{\mathrm{ref}}\succ0\), and the largest generalized eigenvalue is exactly \(D\). All rewards lie in \([0,1]\) and all features have norm at most one. Thus the whole sign family lies in the same exact shell.

It remains to lower-bound its risk. At a hard context, the two target action laws associated with opposite signs are
\[
s_+=(\alpha,\alpha e^{-2t},1-\alpha-\alpha e^{-2t}),
\qquad
s_-=(\alpha e^{-2t},\alpha,1-\alpha-\alpha e^{-2t}).
\]
For any action law \(u\), direct expansion of the two Kullback--Leibler divergences gives
\[
\inf_u\{\operatorname{KL}(u\Vert s_+)+\operatorname{KL}(u\Vert s_-)\}
=-2\log\sum_a\sqrt{s_+(a)s_-(a)}.
\]
Indeed, after normalizing the coordinatewise geometric mean of \(s_+\) and \(s_-\), the left side equals twice a Kullback--Leibler divergence minus twice the logarithm of its normalizer. Here the affinity is
\[
\sum_a\sqrt{s_+(a)s_-(a)}=1-\alpha(1-e^{-t})^2.
\]
For \(0\le t\le1\), the elementary inequalities \(1-e^{-t}\ge t/2\) and \(-\log(1-y)\ge y\) therefore give the conditional two-regret separation
\[
\inf_u\eta^{-1}\{\operatorname{KL}(u\Vert s_+)+\operatorname{KL}(u\Vert s_-)\}
\ge\frac{\alpha t^2}{2\eta}.
\]

Neighboring sign vectors differ only in the two Bernoulli means at one hard context. For \(r,s\in[m,1-m]\), the inequality \(\log x\le x-1\) gives
\[
\operatorname{kl}(r\Vert s)\le\frac{(r-s)^2}{s(1-s)}.
\]
Consequently the one-observation Kullback--Leibler divergence between neighboring laws is at most
\[
\frac{8p\gamma^2}{Hk\,m(1-m)}.
\]
Product additivity gives the same expression multiplied by \(n\). The elementary Pinsker inequality \(\operatorname{TV}(Q,R)^2\le\operatorname{KL}(Q\Vert R)/2\), obtained by applying \((1+y)\log(1+y)+(1-y)\log(1-y)\ge y^2\) after coarsening to the total-variation maximizing event, shows that every neighboring pair has total variation at most \(1/2\) whenever
\[
n\le c_1(C,D,\eta)\frac{k}{p\gamma^2}.
\]

For completeness, if two nonnegative losses \(L_+,L_-\) obey \(L_+(a)+L_-(a)\ge s\) pointwise, integration against the smaller of two data densities gives
\[
\mathbb E_+L_++\mathbb E_-L_-
\ge s\{1-\operatorname{TV}(P_+,P_-)\}.
\]
Apply this inequality to every neighboring pair of sign vertices and average over the uniform prior on \(\{-1,1\}^k\). Each hard context has probability \(1/(Hk)\), and the hard contexts have total probability \(1/H\). Under the preceding testing condition, the Bayes regularized regret is therefore at least
\[
\frac{1}{4H}\frac{\alpha t^2}{2\eta}
=\frac{\alpha t^2}{8H\eta}.
\]
Choose
\[
t^2=16H\eta\epsilon/\alpha.
\]
There is a positive \(\epsilon_0(C,D,\eta)\) for which this choice satisfies \(t\le1\) and \(\gamma=t/\eta\le\gamma_0\) whenever \(0<\epsilon\le\epsilon_0(C,D,\eta)\). The Bayes regret is then at least \(2\epsilon\). Moreover,
\[
p\gamma^2=\frac{\alpha}{D}\frac{t^2}{\eta^2}
=\frac{16H\epsilon}{D\eta}.
\]
Since \(k\ge d/4\), the testing condition contains
\[
n\le c_2(C,D,\eta)\frac{dD\eta}{\epsilon}.
\]
The minimax risk over the full exact shell is at least the Bayes risk over this finite subfamily. Thus it exceeds \(\epsilon\) throughout the last sample-size range. Decreasing \(c_2\) and \(\epsilon_0\) if necessary to absorb integer rounding, \(\text{def:sample-complexity}\) gives the claimed lower bound. The only restriction used was \(D<E=e^\eta\), not \(2D-1<E\).